\section{Why compress actions at all?}

Language arrives as a sequence of symbols with substantial structure. Robot
actions usually arrive as densely sampled numerical signals. Adjacent timesteps
often contain redundant information, but treating every scalar as a separate
token asks a transformer to generate that redundancy one decision at a time.

Consider a one-second chunk with 30 timesteps and seven action dimensions. With
one token per scalar, that is 210 action tokens before any delimiters. A
representation using 24 tokens would require far fewer sequential predictions
for the same physical horizon. These numbers are illustrative, not experimental
results.

Shorter sequences offer fewer opportunities to depart from an expert prefix.
They also reduce the number of sequential decoding steps. Neither benefit
guarantees a better policy: compression can discard useful control detail or
make the remaining tokens harder to predict.

Downsampling is one possible way to reduce redundancy, but moving from 30 Hz to
15 Hz is not automatically harmless. Whether it preserves behavior depends on
the controller, dynamics, and task. A learned tokenizer should instead be judged
on which distinctions it preserves at a given token budget.

The target is therefore not the shortest possible sequence. It is a compact
sequence that preserves task-relevant motion and remains predictable from
observations, language, and preceding tokens.
